\section{The six-sector kernel and the triple coefficient theorem}\label{tsc:triple}
Let $\mathbf s=(s_1,s_2,s_3)$ and let $\mathbf a=(a_1,a_2,a_3)$ have
pairwise distinct entries modulo one. For each $k\in\{1,2,3\}$, let $i,j$
be the other two indices in increasing order and put
\begin{equation}\label{tsc:colors}
 z_{ik}=e(a_k-a_i),\qquad z_{jk}=e(a_k-a_j).
\end{equation}
Define
\begin{align}\label{tsc:kernel}
 \calT(\mathbf s;\mathbf a)=\sum_{k=1}^3\bigg[&
  e^{-\pi\ii(s_i+s_j-s_k)/2}
       T(s_i,s_j,s_k;z_{ik},z_{jk})\notag\\
 &+e^{\pi\ii(s_i+s_j-s_k)/2}
       T(s_i,s_j,s_k;z_{ik}^{-1},z_{jk}^{-1})\bigg].
\end{align}
All six terms are entire in $\mathbf s$ by Theorem~\ref{tsc:T-entire}.
For real spectral variables the terms in each bracket are conjugate, but
\eqref{tsc:kernel} itself is the complex-analytic definition.

\begin{theorem}[Six-sector product identity]\label{tsc:master}
For all $\mathbf s\in\CC^3$,
\begin{equation}\label{tsc:master-eq}
 \left\langle\prod_{i=1}^3\calB_{s_i}(x-a_i),1\right\rangle
       =(2\pi)^{-s_1-s_2-s_3}\calT(\mathbf s;\mathbf a).
\end{equation}
The left side is the separated distribution product. For $\Re s_i>0$
it is an ordinary integral. In particular,
\begin{equation}\label{tsc:T000-check}
 \calT(0,0,0;\mathbf a)=2.
\end{equation}
\end{theorem}
\begin{proof}
First take $\Re s_i>1$, so all three Fourier series are absolutely
convergent. Integration leaves exactly the triples of nonzero integers
$(n_1,n_2,n_3)$ with $n_1+n_2+n_3=0$. Such a triple has either two
positive entries and one negative entry, or the reverse. If $k$ is the
negative singleton, write $n_i=m$, $n_j=n$, $n_k=-(m+n)$ with $m,n\ge1$.
The Fourier powers contribute
\[
 (2\pi)^{-s_1-s_2-s_3}
   e^{-\pi\ii(s_i+s_j-s_k)/2}
          m^{-s_i}n^{-s_j}(m+n)^{-s_k}.
\]
The translation factors contribute
$e(m(a_k-a_i)+n(a_k-a_j))=z_{ik}^m z_{jk}^n$.
Reversing all three signs produces the other term in the bracket. The
six disjoint sectors exhaust the Fourier constraint and give
\eqref{tsc:master-eq}.

Both sides are entire in all three variables. For the left side this was
proved locally at separated singularities; for the right side it follows
from the colored Mellin theorem. The identity theorem, applied successively
to the three variables, extends the equality to all of $\CC^3$.
For $\Re s_i>0$, the only singularity near $a_i$ is an integrable
power $t^{s_i-1}$, so the stated ordinary-integral interpretation holds.
Finally, separated delta functions have zero pairwise products, and
\[
 (\delta_{a_1}-1)(\delta_{a_2}-1)(\delta_{a_3}-1)
    =\delta_{a_1}+\delta_{a_2}+\delta_{a_3}-1.
\]
Pairing with one gives \eqref{tsc:T000-check}.
\end{proof}

\subsection{All Stieltjes indices}
Set
\begin{equation}\label{tsc:E}
 A=\gamma+\log(2\pi),\qquad
 E(u)=\Gamma(1+u)(2\pi)^{-u}
 =\exp\left(-Au+\sum_{k\ge2}\frac{(-1)^k\zeta(k)}k u^k\right).
\end{equation}
The power series is used near zero. Its coefficients are explicit Bell
polynomials in $A,\zeta(2),\ldots$.

\begin{theorem}[All-index triple Stieltjes closure]\label{tsc:stieltjes-triple}
For every nonnegative triple $\mathbf m$ and every separated real shift
triple $\mathbf a$,
\begin{equation}\label{tsc:all-index}
 \boxed{\displaystyle
 \calJ_{\mathbf m}(\mathbf a)=
  \left(\prod_{i=1}^3m_i!\right)
  [u_1^{m_1+1}u_2^{m_2+1}u_3^{m_3+1}]
       \left(\prod_{i=1}^3E(u_i)\right)\calT(\mathbf u;\mathbf a).}
\end{equation}
Every coefficient on the right is a finite combination of jets of the six
entire colored Tornheim functions. No diagonal Laurent finite part or
unspecified order of limits is present.
\end{theorem}
\begin{proof}
Multiply \eqref{tsc:master-eq} by $\prod_i\Gamma(u_i)$ and use
\eqref{tsc:Zlaurent}. The coefficient of
$u_1^{m_1}u_2^{m_2}u_3^{m_3}$ on the left is
$\calJ_{\mathbf m}/\prod_i m_i!$. Indeed, a residue factor in coordinate
$i$ has exponent $-1$ in that coordinate and no dependence on $u_i$ in
any other factor; it cannot contribute to a nonnegative exponent there.
On the right, $\Gamma(u_i)=\Gamma(1+u_i)/u_i$. Moving the three simple
coordinate denominators into the coefficient extraction gives
\eqref{tsc:all-index}. This also proves directly that simultaneous or
iterated coefficient extraction has the same meaning here.
\end{proof}

\subsection{The fully regularized three-digamma formula}
Define $\mathscr D_A=\prod_{i=1}^3(\partial_{s_i}-A)$. Since each
$\gamma_0=-\psi$, the first case of \eqref{tsc:all-index} is
\begin{equation}\label{tsc:digamma-operator}
 \FP\int_0^1\prod_{i=1}^3\psi(\{x-a_i\})\dd x
       =-\left.\mathscr D_A\calT(\mathbf s;\mathbf a)\right|_{\mathbf s=0}.
\end{equation}
This is already a finite, fully normalized answer to the first concrete
target in \cite{tsc:shj}. An equivalent three-term real form is useful.
Put $\alpha=A+\pi\ii/2$. Then
\begin{align}\label{tsc:digamma-three}
 \FP\int_0^1\prod_{i=1}^3\psi(\{x-a_i\})\dd x
 =-2\Re\sum_{k=1}^3
 \left.(\partial_r-\alpha)(\partial_s-\alpha)
       (\partial_t-\overline\alpha)
       T(r,s,t;z_{ik},z_{jk})\right|_{r=s=t=0}.
\end{align}
Only derivatives of order at most one in each of the three spectral
variables occur. The next section turns these derivatives into ordinary
convergent integrals.

\subsection{Translation and spectral lowering}
The Fourier definition gives $\partial_x\calB_s=\calB_{s-1}$ as an
entire distributional identity. Differentiating the separated product yields
\begin{equation}\label{tsc:lowering}
 \partial_{a_j}\calT(\mathbf s;\mathbf a)
       =-2\pi\calT(\mathbf s-\mathbf e_j;\mathbf a),
 \qquad
 \sum_{j=1}^3\calT(\mathbf s-\mathbf e_j;\mathbf a)=0.
\end{equation}
The second equality also follows from simultaneous translation invariance.
These are identities for the complete six-sector combination, not a license
to discard sectors or their phases.
